\chapter{An example of policies' comparison}
\label{ch:policy}
The last version of the model shows a stable behaviour and enough heterogeneity and distributional inequalities. It makes somehow suitable to sketch an answer to the original question: what will happen if an UBI will be introduced?

A simple experiment has been set up, creating five scenarios and five repetition for each of them. The scenarios are:
\begin{enumerate}
    \item the baseline model from the previous chapter.
    \item a model with increased profit taxation.
    \item a model with the introduction of a wealth tax paid at the end of each period.
    \item a model with an UBI, implemented as if the unemployment benefits from the previous models are paid regardless the employment status.
    \item a model with an UBI and a wealth tax.
\end{enumerate}

\begin{figure}
    \centering
    \begin{subfigure}[t]{0.49\textwidth}
        \centering
        \includegraphics[width=\textwidth]{pol_i}
        \caption{Dynamics of inflation after the first few periods}
        \label{fig:pol_i}
    \end{subfigure}%
    ~
    \begin{subfigure}[t]{0.49\textwidth}
        \centering
        \includegraphics[width=\textwidth]{pol_u}
        \caption{Dynamics of unemployment after the first few periods}
        \label{fig:pol_u}
    \end{subfigure}
    \begin{subfigure}[t]{0.49\textwidth}
        \centering
        \includegraphics[width=\textwidth]{pol_gdp}
        \caption{Dynamics of nominal GDP}
        \label{fig:pol_sb}
    \end{subfigure}%
    ~
    \begin{subfigure}[t]{0.49\textwidth}
        \centering
        \includegraphics[width=\textwidth]{pol_g}
        \caption{Dynamics of the share of public consumption}
        \label{fig:pol_wd}
    \end{subfigure}
    \begin{subfigure}[t]{0.49\textwidth}
        \centering
        \includegraphics[width=\textwidth]{pol_wgini}
        \caption{Dynamics of Gini index calculated on wages}
        \label{fig:pol_sdist}
    \end{subfigure}%
    ~
    \begin{subfigure}[t]{0.49\textwidth}
        \centering
        \includegraphics[width=\textwidth]{pol_mgini}
        \caption{Dynamics of Gini index calculated on wealth after the first few periods}
        \label{fig:pol_bdist}
    \end{subfigure}
    \caption{Dynamics of the AB model with innovation}
\end{figure}

The emerging results are not necessarily intuitive.
The introduction of an UBI accelerates the nominal growth of the model, its instability and move consumption from the public to the private sector.
But at the same time it is the policy that most increase the Gini index, appearing as the policy less able to reduce inequalities.
This is reasonably because such a policy sustains the demand through a direct transfer of wealth from the public to the private sector, reducing the relative government consumption (which is not included in the Gini calculation) and increasing the private one.

Old Keynesian policies based on taxation slow down the economy but at the same time are very efficient in reducing inequalities.
In fact they both are targeted towards the wealthiest households (one targeting the flows, the other the stock) aiming to reduce their consumption and increasing the public one (which again is invisible to Gini calculation).
Those reveal themselves as effective deflationary policies, as stated by some post-Keynesian theories.

Mixing the two policies appears to be in this model\footnote{The hidden assumption is that this model is a reliable representation of the reality, which I would argue against.} an interesting proposal: the version which mix an UBI and a wealth tax shows faster growth than the baseline model but not higher inequality.
The taxation neutralizes the transfer of wealth towards richer households, making the UBI a purely redistributive policy which sustain the aggregate private demand.

This is clearly a too simple experiment in a toy-model, but it shows very clearly and effectively the path and the strategies a researcher should follow to answer such a question as the original one.

Finally, in the teaching spirit of the work it would be very interesting to transform the model into a video game, giving the player agency in changing some parameters of the model, especially those related to the public sector, to explore the dynamics of the model as a playground to observe some Economics laws.